\section{USAMO 2016}
        \begin{enumerate}
        	\item (\textit{USAMO 2016}) Let $X_1, X_2, \ldots, X_{100}$ be a sequence of mutually distinct nonempty subsets of a set $S$. Any two sets $X_i$ and $X_{i+1}$ are disjoint and their union is not the whole set $S$, that is, $X_i\cap X_{i+1}=\emptyset$ and $X_i\cup X_{i+1}\neq S$, for all $i\in\{1, \ldots, 99\}$. Find the smallest possible number of elements in $S$.


 

	\item (\textit{USAMO 2016}) Prove that for any positive integer $k,$

 \[\left(k^2\right)!\cdot\prod_{j=0}^{k-1}\frac{j!}{\left(j+k\right)!}\]
is an integer.


 

	\item (\textit{USAMO 2016}) Let $\triangle ABC$ be an acute triangle, and let $I_B, I_C,$ and $O$ denote its $B$-excenter, $C$-excenter, and circumcenter, respectively. Points $E$ and $Y$ are selected on $\overline{AC}$ such that $\angle ABY = \angle CBY$ and $\overline{BE}\perp\overline{AC}.$ Similarly, points $F$ and $Z$ are selected on $\overline{AB}$ such that $\angle ACZ = \angle BCZ$ and $\overline{CF}\perp\overline{AB}.$


 Lines $\overleftrightarrow{I_B F}$ and $\overleftrightarrow{I_C E}$ meet at $P.$ Prove that $\overline{PO}$ and $\overline{YZ}$ are perpendicular.


 

	\item (\textit{USAMO 2016}) Find all functions $f:\mathbb{R}\rightarrow \mathbb{R}$ such that for all real numbers $x$ and $y$, 
 \[(f(x)+xy)\cdot f(x-3y)+(f(y)+xy)\cdot f(3x-y)=(f(x+y))^2.\]


 

	\item (\textit{USAMO 2016}) An equilateral pentagon $AMNPQ$ is inscribed in triangle $ABC$ such that $M\in\overline{AB},$ $Q\in\overline{AC},$ and $N, P\in\overline{BC}.$ Let $S$ be the intersection of $\overleftrightarrow{MN}$ and $\overleftrightarrow{PQ}.$ Denote by $\ell$ the angle bisector of $\angle MSQ.$


 Prove that $\overline{OI}$ is parallel to $\ell,$ where $O$ is the circumcenter of triangle $ABC,$ and $I$ is the incenter of triangle $ABC.$


 

	\item (\textit{USAMO 2016}) Integers $n$ and $k$ are given, with $n\ge k\ge 2.$ You play the following game against an evil wizard.


 The wizard has $2n$ cards; for each $i = 1, ..., n,$ there are two cards labeled $i.$ Initially, the wizard places all cards face down in a row, in unknown order.


 You may repeatedly make moves of the following form: you point to any $k$ of the cards. The wizard then turns those cards face up. If any two of the cards match, the game is over and you win. Otherwise, you must look away, while the wizard arbitrarily permutes the $k$ chosen cards and turns them back face-down. Then, it is your turn again.


 We say this game is \textit{winnable} if there exist some positive integer $m$ and some strategy that is guaranteed to win in at most $m$ moves, no matter how the wizard responds.


 For which values of $n$ and $k$ is the game winnable?


 

\end{enumerate}